% ===================================================================
%  TAVOLA N. 14 — La strada. --- I negozianti.
%  Traduction 2026 d'apres texto/io, controlee sur texto/fr et
%  texto/en. Consigne : voir texto/it/10-tavola-01.tex.
%
%  LES DEUX CARACTERES DE TOUS LES LIVRES DU MONDE PORTENT LE NOM DE
%  CE PAYS ET DE SA CAPITALE, ET L'ALINEA 17 EST UNE IMPRIMERIE.
%
%  Le romain et l'italique : « roman » et « italic » en anglais,
%  « Antiqua » et « Kursive » en allemand, « romain » et « italique »
%  en francais. L'un est le caractere que les humanistes de Rome ont
%  taille sur l'ecriture des inscriptions antiques ; l'autre celui
%  qu'Alde Manuce a fait graver a Venise en 1501 sur la cursive des
%  chancelleries italiennes, pour ses editions de poche.
%
%  Or l'italien ne les appelle pas ainsi. Il dit tondo, le
%  rond, et corsivo, le courant. UN PAYS NE NOMME PAS SES
%  PROPRES CARACTERES PAR SON NOM : c'est de l'exterieur qu'on voit
%  d'ou vient une chose. Les trente-trois autres colonnes composent
%  leurs notes en « italique » sans y penser ; celle-ci les compose en
%  « corsivo », et c'est la seule ou le mot ne dise rien d'un lieu.
%
%  Le tableau 14 allemand avait trouve que Setzer, Drucker et
%  Buchbinder etaient partis de Mayence vers 1450 dans toute l'Europe
%  de l'est. Voici l'autre moitie du meme mouvement : L'ALLEMAGNE A
%  EXPORTE LES NOMS DES OUVRIERS, L'ITALIE LES FORMES DES LETTRES.
%
%  ET LA BOTTEGA DE CETTE PAGE A TROIS SOEURS QUI NE LUI RESSEMBLENT
%  PLUS.
%
%  Le magasin s'appelle en italien bottega, du grec
%  « apothēkē » par le latin « apotheca » : LE MAGASIN, l'endroit ou
%  l'on serre. Le meme mot est devenu :
%     en allemand   « Apotheke »  la pharmacie — c'est celle du
%                   tableau 4, alinea 6 de la seconde scene ;
%     en francais   « boutique »  le petit commerce ;
%     en espagnol   « bodega »    la cave a vin ;
%     en italien    « bottega »   l'atelier et le magasin.
%  UN SEUL MOT GREC, QUATRE LANGUES, QUATRE RETRECISSEMENTS
%  DIFFERENTS, et chacune croit tenir le sens propre. Aucune ne l'a :
%  le grec disait simplement « ce qu'on met de cote ».
%
%  Le reste de la rue se nomme par le suffixe « -aio » ou « -iere »,
%  qui n'est que le latin « -arius » : orologiaio, cappellaio,
%  guantaio, libraio, tabaccaio, pasticciere, gioielliere, barbiere.
%  L'allemand composait deux racines, l'italien colle un suffixe. LES
%  DEUX FONT LE MEME TRAVAIL, ET AUCUNE DES DEUX N'A BESOIN
%  D'EMPRUNTER.
% ===================================================================

%%K t14-apar-1 apar
\VUsaut{4.0mm}
\VUtitre{15.0pt}{60}{TAVOLA N. 14}
\VUsaut{3.0mm}

%%K t14-tit-1 sub
\VUpk{11.4pt}{40}{La strada. --- I negozianti.}
\VUsaut{3.0mm}

%%K t14-01-1 p
1. --- Il mio amico Giovanni, che abita in campagna, è venuto a passare tre giorni
dalla mia famiglia. \VUgras{Fino} a oggi non aveva mai avuto occasione di vedere una
grande città, e così passiamo tutte le giornate a spasso e io gli spiego quello che
non conosce.

%%K t14-02-1 p
2. --- Prima siamo andati a vedere l'\VUgras{osservatorio} \textsuperscript{(1)}, che
gli è piaciuto moltissimo. Tornando per la \VUgras{via} \textsuperscript{(3)} dove ci
sono i \VUgras{bagni turchi} \textsuperscript{(2)}, abbiamo incontrato un
\VUgras{funerale} \textsuperscript{(4)}. In testa al \VUgras{corteo funebre}
camminava il \VUgras{clero}, davanti al \VUgras{carro funebre} in cui era stata messa
la \VUgras{bara}. Molti amici seguivano la famiglia in \VUgras{lutto}.

%%K t14-02-2 p
Passato il corteo, abbiamo attraversato la strada davanti alla grande
\VUgras{birreria} \textsuperscript{(5)}, dove molti \VUgras{clienti} bevevano
\VUgras{birra} sulla terrazza, riparati dalla \VUgras{tettoia} \textsuperscript{(6)} a
vetri.

%%K t14-03-1 p
3. --- Giovanni non capiva lo \VUgras{stemma} messo fra il negozio del
\VUgras{vinaio} \textsuperscript{(7)} e quello del \VUgras{negoziante di musica}
\textsuperscript{(10)}. Mi ha chiesto che cosa volesse dire; gli ho detto che indicava
il \VUgras{consolato} \textsuperscript{(8)} inglese, e gli ho mostrato il
\VUgras{console} \textsuperscript{(9)} in piedi sul balcone.

%%K t14-tit-2 sub
\VUpk{10.2pt}{40}{Guardare le vetrine.}
\VUsaut{3.0mm}

%%K t14-04-1 p
4. --- Naturalmente ci siamo fermati davanti all'esposizione di
\VUgras{strumenti musicali}, \VUgras{armonium}, \VUgras{organi} e
\VUgras{pianoforti}; abbiamo guardato le copertine illustrate degli \VUgras{spartiti}
e dei \VUgras{pezzi} per pianoforte, e abbiamo ammirato la \VUgras{lira} di legno
dorato che serve da insegna.

%%K t14-05-1 p
5. --- Le nuvole di polvere alzate dagli \VUgras{spazzini} \textsuperscript{(11)} e
dalla \VUgras{spazzatrice} \textsuperscript{(12)} ci hanno costretti a passare in
fretta davanti alle belle \VUgras{camere complete} del \VUgras{mobiliere}
\textsuperscript{(13)} e all'esposizione di \VUgras{stufe} e di \VUgras{lampade} del
\VUgras{ferramenta} \textsuperscript{(14)}.

%%K t14-06-1 p
6. --- In quel punto abbiamo dovuto davvero lasciare il marciapiede, perché era
ingombro di pacchi che venivano scaricati da un \VUgras{furgone da trasloco}
\textsuperscript{(16)} per essere portati in un \VUgras{negozio}
\textsuperscript{(15)} appena affittato.

%%K t14-06-2 p
Così ci sono sfuggite le belle carrozze del \VUgras{carrozziere}
\textsuperscript{(17)}, ma non ci siamo persi l'\VUgras{imballatore}
\textsuperscript{(19)} che fabbricava una cassa per il suo vicino, il
\VUgras{negoziante di biciclette e automobili} \textsuperscript{(18)}. Poi, siccome era
già piuttosto tardi, siamo tornati a casa.

%%K t14-tit-3 sub
\VUpk{10.2pt}{40}{Un giro per la città.}
\VUsaut{3.0mm}

%%K t14-07-1 p
7. --- Il giorno dopo mia madre ha proposto di far fare una fotografia a Giovanni e mi
ha chiesto di accompagnarlo dal \VUgras{fotografo} \textsuperscript{(20)}. Dopo pranzo
siamo andati allo \VUgras{studio} dell'artista, dove abbiamo dovuto aspettare
parecchio. Per passare il tempo abbiamo guardato i \VUgras{ritratti}
\textsuperscript{(22)} appesi al muro.

%%K t14-07-2 p
Quando è venuto il nostro turno è stata cosa breve, e appena il mio amico ha finito di
posare davanti alla \VUgras{macchina} \textsuperscript{(21)}, siamo scesi.

%%K t14-08-1 p
8. --- Al primo piano abbiamo visto, dalla porta aperta del \VUgras{barbiere}
\textsuperscript{(23)}, il suo garzone che tagliava i capelli a un cliente.

%%K t14-08-2 p
Tornati in strada, siamo passati davanti alla \VUgras{tabaccheria}
\textsuperscript{(24)}. Giovanni si è divertito tanto per la \VUgras{lanterna}
\textsuperscript{(25)} rossa e la \VUgras{grande pipa} che serve da \VUgras{insegna}
\textsuperscript{(26)} che è andato a sbattere contro due vecchi signori che
chiacchieravano davanti a una \VUgras{presa di tabacco} \textsuperscript{(27)}.

%%K t14-tit-4 sub
\VUpk{10.2pt}{40}{La pasticceria.}
\VUsaut{3.0mm}

%%K t14-09-1 p
9. --- Le \VUgras{banconote} e le \VUgras{monete} di tutti i paesi esposte nella
\VUgras{vetrina} del \VUgras{cambiavalute} \textsuperscript{(29)} ci hanno trattenuti
un momento, ma siccome era l'ora della merenda siamo entrati senza altro indugio nella
\VUgras{pasticceria} \textsuperscript{(31)}.

%%K t14-10-1 p
10. --- Davanti a tutte le \VUgras{cose buone} del negozio --- \VUgras{torte},
\VUgras{panpepato}, \VUgras{biscotti} e \VUgras{dolciumi} di ogni genere --- non
riuscivamo a scegliere. Alla fine Giovanni ha preso i \VUgras{bignè alla crema}, e io
ho preso solo una piccola \VUgras{crostatina di ciliegie}, perché il dolce
\VUgras{mi fa venire il mal di denti} e preferisco non andare troppo spesso dal
\VUgras{dentista} \textsuperscript{(30)}.

%%K t14-tit-5 sub
\VUpk{10.2pt}{40}{Scrivere a macchina.}
\VUsaut{3.0mm}

%%K t14-11-1 p
11. --- Per far vedere al mio amico una \VUgras{macchina da scrivere}, di cui aveva
sentito parlare ma che non aveva mai visto, siamo entrati nell'\VUgras{ufficio}
\textsuperscript{(32)} di mio \VUgras{cugino} \textsuperscript{(34)}. L'abbiamo trovato
che dettava le sue lettere alla sua \VUgras{segretaria} \textsuperscript{(35)}, che è
\VUgras{stenografa}. Appena una lettera era finita, una \VUgras{dattilografa}
\textsuperscript{(33)} la ricopiava a macchina. Non volevamo disturbare il lavoro di
mio parente e siamo rimasti solo pochi minuti.

%%K t14-12-1 p
12. --- Sotto l'ufficio di mio cugino abita un grande
\VUgras{orologiaio e gioielliere} \textsuperscript{(37)}, che è anche argentiere. Il
suo splendido negozio ha molti \VUgras{orologi}, \VUgras{pendole},
\VUgras{orologi da tasca}, \VUgras{catene} e \VUgras{gioielli} preziosi come
\VUgras{collane di perle}, \VUgras{braccialetti} d'oro, \VUgras{orecchini} e
\VUgras{anelli} con \VUgras{pietre preziose} e \VUgras{diamanti}. Una delle vetrine è
dedicata tutta all'\VUgras{argenteria} e contiene una grande collezione di
\VUgras{posate} d'argento.

%%K t14-13-1 p
13. --- Girato l'angolo della strada, mi sono accorto che la \VUgras{targa}
\textsuperscript{(36)} accanto al \VUgras{numero} della casa era stata cambiata. Stavo
per dirlo ad alta voce quando abbiamo sentito il suono allegro di una banda militare.
Ci siamo messi a correre verso il reggimento, senza pensare che sarebbe passato davanti
ai negozi del \VUgras{cappellaio} \textsuperscript{(39)} e del \VUgras{guantaio}
\textsuperscript{(40)}, e che l'avremmo visto sfilare altrettanto bene restando davanti
alla vetrina dell'\VUgras{ottico} \textsuperscript{(38)}, piena di \VUgras{occhialini},
\VUgras{occhiali} e \VUgras{binocoli da teatro}.

%%K t14-tit-6 sub
\VUpk{10.2pt}{40}{La sfilata.}
\VUsaut{3.0mm}

%%K t14-14-1 p
14. --- In testa al \VUgras{reggimento} \textsuperscript{(41)} camminava il
\VUgras{tamburo maggiore} \textsuperscript{(42)}, seguito dai \VUgras{tamburini}
\textsuperscript{(43)}, dai \VUgras{trombettieri} \textsuperscript{(44)} e da tutta la
\VUgras{banda}. Il \VUgras{generale} \textsuperscript{(45)}, che era sulla piazza con
il suo aiutante di campo, si era fermato accanto allo \VUgras{zampillo}
\textsuperscript{(49)} della \VUgras{fontana} \textsuperscript{(48)} per guardare la
\VUgras{parata}.

%%K t14-14-2 p
\VUgras{Gli ufficiali di stato maggiore} \textsuperscript{(46)}, il \VUgras{colonnello}
e il \VUgras{tenente colonnello} lo salutavano con la spada. \VUgras{I maggiori} erano
in testa ai loro \VUgras{battaglioni} \textsuperscript{(47)} e ogni \VUgras{capitano}
comandava la sua \VUgras{compagnia}, mentre i \VUgras{tenenti}, i
\VUgras{sottotenenti} e i \VUgras{sottufficiali} stavano al loro posto abituale
accanto ai loro uomini.

%%K t14-15-1 p
15. --- Molti passanti si erano radunati all'\VUgras{incrocio}
\textsuperscript{(50)}; i negozianti --- il \VUgras{venditore di ombrelli}
\textsuperscript{(51)}, il \VUgras{tintore} \textsuperscript{(53)}, l'\VUgras{armaiolo}
\textsuperscript{(54)} --- erano usciti a vedere i \VUgras{soldati}. Tutti i veicoli ---
\VUgras{carrozze da nolo} \textsuperscript{(52)}, \VUgras{carrozze private}
\textsuperscript{(57)} e anche l'\VUgras{annaffiatrice} \textsuperscript{(56)} --- si
erano accostati al marciapiede.

%%K t14-tit-7 sub
\VUpk{10.2pt}{40}{Scene di strada.}
\VUsaut{3.0mm}

%%K t14-15-2 p
Un \VUgras{commesso viaggiatore} \textsuperscript{(55)} con le sue
\VUgras{valigie di campioni} in mano attraversava svelto la strada, e le persone sedute
sull'\VUgras{imperiale} \textsuperscript{(59)} dell'\VUgras{omnibus}
\textsuperscript{(58)} si voltavano per godersi lo spettacolo. Un povero
\VUgras{cantastorie} \textsuperscript{(60)} con il suo \VUgras{organetto}
\textsuperscript{(61)} ha dovuto interrompere la sua \VUgras{canzone}, perché il rumore
dei tamburi copriva la sua voce.

%%K t14-16-1 p
16. --- Quando il reggimento è arrivato davanti al \VUgras{negozio di stoffe}
\textsuperscript{(64)}, le \VUgras{sarte} \textsuperscript{(63)} e i \VUgras{sarti}
\textsuperscript{(62)} hanno lasciato le loro \VUgras{macchine da cucire} e il loro
lavoro per affacciarsi alla finestra. \VUgras{Il negoziante} \textsuperscript{(65)},
che aveva appena messo \VUgras{abiti} nuovi \textsuperscript{(66)} nella sua
\VUgras{vetrina}, ha dovuto spostare i rotoli di \VUgras{panno}
\textsuperscript{(67)} e di \VUgras{tela} esposti sul marciapiede, vicino alla
\VUgras{bocca di scarico} \textsuperscript{(68)}, per paura che la folla li rovesciasse;
perfino un vecchio \VUgras{merciaio ambulante} \textsuperscript{(69)}, con cui una
portinaia stava contrattando delle calze, ha dovuto entrare in un portone per finire la
vendita.

%%K t14-16-2 p
Abbiamo accompagnato il reggimento fino alla caserma \VUgras{e}, molto contenti della
nostra giornata, siamo tornati a casa per cena.

%%K t14-tit-8 sub
\VUpk{10.2pt}{40}{Mestieri e professioni.}
\VUsaut{3.0mm}

%%K t14-17-1 p
17. --- Il giorno dopo mio padre ci ha proposto di portarci a vedere la tipografia del
giornale locale. Abbiamo accettato con entusiasmo.

%%K t14-17-2 p
Il \VUgras{tipografo} è anche \VUgras{libraio} \textsuperscript{(70)},
\VUgras{editore}, \VUgras{cartolaio} e \VUgras{legatore}. Al primo piano ha messo la
\VUgras{redazione} \textsuperscript{(71)} del giornale, e anche il
\VUgras{laboratorio di incisione}, dove lavorano i \VUgras{litografi}
\textsuperscript{(72)} e gli \VUgras{incisori in rame} \textsuperscript{(73)}, e
infine la sala di composizione, dove stanno i \VUgras{compositori}
\textsuperscript{(74)}. La vera e propria \VUgras{sala delle macchine}
\textsuperscript{(75)}, i \VUgras{torchi} e le varie macchine sono al pianterreno.

%%K t14-tit-9 sub
\VUpk{10.2pt}{40}{Giovanni fa moltissime domande.}
\VUsaut{3.0mm}

%%K t14-18-1 p
18. --- Uscendo dalla tipografia, Giovanni si è fermato tutto stupito davanti a una
cassetta di vetro montata su un palo di fronte alla \VUgras{merceria}
\textsuperscript{(77)}, e ha chiesto a che cosa servisse. Mio padre gli ha spiegato che
era un \VUgras{avvisatore d'incendio} \textsuperscript{(76)}. Del resto in città tutto
era una meraviglia per il mio amico, dalla semplice \VUgras{fontanella}
\textsuperscript{(78)} e dal leggero \VUgras{triciclo da consegna}
\textsuperscript{(80)} guidato da un \VUgras{fattorino} \textsuperscript{(79)} in
livrea, fino al \VUgras{furgone postale} \textsuperscript{(81)}.

%%K t14-19-1 p
19. --- Mentre mio padre ci lasciava un momento per parlare con l'\VUgras{esattore}
\textsuperscript{(83)}, che si era fermato a discorrere con un \VUgras{avvocato}
\textsuperscript{(82)} e un \VUgras{notaio} \textsuperscript{(84)}, ci siamo divertiti a
seguire il movimento della strada. Ci passavano davanti le persone più diverse: un
\VUgras{garzone di pasticceria} \textsuperscript{(85)}, che portava in un cesto un
magnifico \VUgras{pasticcio}, veniva dietro a un \VUgras{rigattiere}
\textsuperscript{(86)}; un \VUgras{lampionaio} \textsuperscript{(87)} parlava con un
\VUgras{gasista} \textsuperscript{(88)}; una \VUgras{suora} \textsuperscript{(89)} che
teneva per mano una piccola orfana tornava al suo convento.

%%K t14-tit-10 sub
\VUpk{10.2pt}{40}{Sulla via di casa.}
\VUsaut{3.0mm}

%%K t14-20-1 p
20. --- Proprio mentre mio padre ci raggiungeva, Giovanni è scoppiato a ridere per le
facce nere di fuliggine e la strana tenuta di due \VUgras{spazzacamini}
\textsuperscript{(91)}.

%%K t14-21-1 p
21. --- Volevo comprare a mia madre una pianta dalla \VUgras{fioraia}
\textsuperscript{(92)}, ma mio padre mi ha fatto notare che sarebbe stata molto pesante
da portare e che era meglio prendere delle \VUgras{arance}. Siamo andati dalla
\VUgras{venditrice} \textsuperscript{(94)}, e quando la \VUgras{governante}
\textsuperscript{(93)}, che ne sceglieva per la sua bambina, ha finito la sua spesa,
siamo stati serviti noi.

%%K t14-22-1 p
22. --- Tornando a casa abbiamo incontrato diverse persone che conosciamo: una vecchia
amica della mia famiglia, appoggiata al braccio della sua \VUgras{dama di compagnia}
\textsuperscript{(95)}, l'\VUgras{ingegnere} \textsuperscript{(96)} del genio civile, e
una mia cugina con la sua \VUgras{bambinaia} \textsuperscript{(98)}.

%%K t14-22-2 p
Poi siamo passati davanti alla \VUgras{modista} \textsuperscript{(97)} di mia madre e
davanti a due \VUgras{frati} \textsuperscript{(99)} che erano forestieri in città e che
hanno chiesto a mio padre la strada per la stazione.
\VUsaut{6.0mm}
\VUfilet{22mm}
